Estimating thresholds and benchmarking decoders for quantum error correction requires sampling very large numbers of
detector outcomes from noisy stabilizer circuits, today mostly with the Stim library on CPUs. We
present a Pauli frame sampler for the AMD Alveo U55C that processes one instruction slot per clock cycle for
\NumLanes{} shots at a time, and we run it on four cards in one host that are treated as independent nodes, as in our
earlier multi FPGA statevector simulator. Three changes remove the bottlenecks of a first single card kernel:
noise masks are produced by a separate pipelined stage from a multiplier free counter based generator, the Pauli frame of up
to \NumQubitsMax{} qubits is updated through a one read, one write pipeline whose hazards the compiler removes by
scheduling, and detectors are folded from a per batch measurement record while the next batch runs. Each card holds \NumInstancesPerCard{} kernel instances in separate super logic regions, and every node, instance
and batch draws its own random stream, so the nodes exchange no data. The kernel is bit exact against a software model in
C simulation, in RTL cosimulation and on every card, and its detector rates and decoded logical error rates agree with
Stim for surface codes up to distance 7. On four cards the sampler delivers \NumFpgaFourSurfFive{} shots per second
for a distance 5 surface code, against \NumCpuOneSurfFive{} for one Stim process and
\NumCpuAllSurfFive{} for \NumCpuMaxProcs{} Stim processes on the host's \NumHostThreads{} hardware threads.
